\documentclass[11pt]{article}
\usepackage[margin=1.15in]{geometry}
\usepackage{amsmath,amssymb}
\usepackage{booktabs}
\usepackage{microtype}
\usepackage[hidelinks]{hyperref}
\usepackage{titlesec}
\usepackage[T1]{fontenc}
\usepackage{lmodern}
\titleformat{\section}{\normalfont\large\bfseries}{\thesection}{0.6em}{}
\titleformat{\subsection}{\normalfont\normalsize\bfseries}{\thesubsection}{0.6em}{}
\setlength{\parskip}{0.35em}

\title{\bfseries Generality Is Compression:\\
An Operational Measure, and a Representation That Fails It}
\author{Shehzad Ahmed\\\small Independent University, Bangladesh}
\date{September 2026}

\begin{document}
\maketitle

\begin{abstract}
\noindent
``General'' is used to mean two different things. One is \emph{breadth}: the
system handles many cases because many cases were stored. The other is
\emph{generality}: the system handles many cases because one structure covers
them. The distinction is standard as intuition---the novice sorts physics
problems by surface features, the expert by governing principle---and is
usually left there.

We give it an operational measure. If generality is compression, then a general
representation should hold description length roughly fixed while the world it
describes grows, and its \emph{coverage per stored item} should therefore rise
with environment size, while a merely broad representation's falls. We test this
on an agent whose memory can be indexed two ways that differ in nothing but the
key: by \emph{location}, or by \emph{local situation}. Across environments of
25, 64, and 144 states, coverage per stored item moves in opposite directions,
$0.99 \to 0.83 \to 0.75$ for location and $2.13 \to 4.12 \to 6.93$ for
situation. The situation-keyed store also crosses \emph{below} the
location-keyed one in size while covering seven times as many cases per item.

The negative half is sharper than the positive. Location-keyed memory does not
merely fail to transfer to unseen environments---it \emph{harms}, at every
scale, and is worse than a store whose contents have been randomly shuffled.
The asymmetry has a mechanism: a shuffled store is harmless because it rarely
matches and therefore rarely fires, whereas the location-keyed store fires
constantly and the world behind its matches has moved. \textbf{It retrieves
reliably and retrieves wrong.} All seven predictions were preregistered and
held on 20{,}000 seeds disjoint from any exploratory range.
\end{abstract}

\section{Two things called general}

A chess engine plays chess. A translator translates. Adding both to a system
makes it handle more cases, and nothing about the second helps with the first.
Call this \textbf{breadth}: coverage purchased by accumulation.

Contrast the classic finding on expertise. Given a set of physics problems,
novices sort them by surface---inclined planes here, pulleys there---while
experts sort the same problems by governing principle, placing a ramp and a
pulley together because both are conservation of energy. The expert's
organisation is smaller and covers more. Call this \textbf{generality}:
coverage purchased by structure.

The distinction is intuitive and rarely measured. This paper offers a measure
and applies it to a case where the two representations differ in exactly one
respect, so that the comparison is clean.

\section{Why compression is the right test}

The reason to expect compression to separate them is not analogy. It follows
from what compression is.

Shannon's entropy $H(P) = \sum_i p_i \log(1/p_i)$ is the minimum average bits
per symbol achievable if one's code matches the world. Coding a world $P$ with
a scheme optimised for $Q$ costs the cross-entropy
$H(P,Q) = \sum_i p_i \log(1/q_i)$, which exceeds $H(P)$ by the
Kullback--Leibler divergence and is minimised exactly when $Q = P$.

That minimisation property is not one convenient choice among many. If one
asks which loss $F$ makes $\sum_i P_i F(Q_i)$ minimal at $Q = P$ subject to
$\sum_i Q_i = 1$, the Lagrange condition gives $P_i F'(Q_i) = \lambda$, hence
$F'(x) \propto 1/x$ and $F = -\log$ uniquely. Cross-entropy is forced.

Two consequences matter here. First, prediction and compression are the same
operation: any predictor $Q$ converts, via arithmetic coding, into a compressor
achieving approximately $\sum -\log Q$ bits, so minimising predictive loss is
minimising code length. Second, description length is the natural formalisation
of ``one structure under many things''---the Kolmogorov intuition that the
shortest program producing a body of data is the structure in it.

Breadth and generality are therefore distinguishable in principle by how
description length behaves as the described world grows. Breadth must grow with
it: each new case is a new entry. Generality need not, because the entries were
never about cases.

\section{An operational measure}

Practical description length is not computable, but a proxy is. For a system
with an explicit store of $N$ items that are retrieved during operation, define

\[
\text{coverage} \;=\; \frac{\text{retrievals per episode}}{N}
\]

the number of situations each stored item is actually used on. A representation
whose items name individual cases has coverage near one and falling: as the
world grows there are more cases, each item still names one, and the chance of
revisiting any particular case declines. A representation whose items name
recurring \emph{situations} has coverage above one and rising: the number of
distinct situations does not grow with the size of the world, so each item is
matched more often.

This yields the prediction we test:

\begin{quote}
\textbf{As environment size grows, coverage per stored item falls for a
case-indexed representation and rises for a structure-indexed one.}
\end{quote}

The prediction is about a \emph{direction of change with scale}, which is
harder to satisfy by accident than a difference at one size.

\section{Experiment}

\subsection{Task and representations}

An agent traverses a grid whose obstacles are re-randomised for every episode,
so route memorisation cannot help. During training it stores experiences and
consults the store when scoring candidate moves. We vary only the key:

\begin{itemize}\itemsep1pt
\item \textbf{Location key:} \texttt{at (x,y) $\to$ move}. Names the cell.
\item \textbf{Situation key:} \texttt{blk<4 bits>|goal(dx,dy) $\to$ move}. Names
which of the four neighbours are blocked, plus the coarse direction of the
goal.
\end{itemize}

The agent already computes neighbour blockage in order to bias toward the goal,
so the situation key uses \emph{no information the location key lacked}. The
two conditions are identical agents differing in one string.

\subsection{Transfer design}

Training runs on four novel environments. The agent then attempts one it has
never seen, three times from the same trained state, differing only in the
store: \textbf{intact}, \textbf{wiped} (empty), or \textbf{shuffled} (same
size, same keys, values permuted). The shuffled arm is the control that
separates a store's \emph{content} from the mechanics of having a store at all.

Goal bias is held constant across arms and the held-out episode uses an
identical random stream, so the three arms are paired and the statistic is the
per-seed difference.

\subsection{Preregistration}

Seven predictions with numeric thresholds---set below the exploratory point
estimates so the test is a replication rather than a restatement---were
committed to version control before the first replication seed was drawn, along
with the analysis plan and a stopping rule of one run. Seeds
$500001$--$520000$ are disjoint from every range used in exploratory work.
The primary outcome was declared to be the intact-versus-\emph{shuffled}
contrast, not intact-versus-wiped: the question is whether content carries
structure, not whether having a store does anything.

\section{Results}

\begin{table}[h]
\centering
\begin{tabular}{llccccc}
\toprule
States & Key & Store & Retrievals & \textbf{Coverage} & intact $-$ wiped & intact $-$ shuffled \\
\midrule
25 & location & 6.77 & 6.7 & \textbf{0.99} & $-0.0108$ & $-0.0126$ \\
25 & situation & 8.51 & 18.1 & \textbf{2.13} & $+0.0425$ & $+0.0309$ \\
64 & location & 11.55 & 9.6 & \textbf{0.83} & $-0.0040$ & $-0.0042$ \\
64 & situation & 12.05 & 49.6 & \textbf{4.12} & $+0.0442$ & $+0.0352$ \\
144 & location & 17.24 & 13.0 & \textbf{0.75} & $-0.0043$ & $-0.0029$ \\
144 & situation & 14.43 & 100.0 & \textbf{6.93} & $+0.0207$ & $+0.0159$ \\
\bottomrule
\end{tabular}
\caption{Held-out transfer, $n = 20{,}000$ paired seeds per cell. All six
differences exclude zero. All seven preregistered predictions held.}
\end{table}

\paragraph{Coverage diverges, as predicted.} Location: $0.99 \to 0.83 \to
0.75$. Situation: $2.13 \to 4.12 \to 6.93$. The environment grew $5.8\times$ in
area; the location store grew $2.6\times$ and the situation store $1.7\times$.
Sub-linear, not flat---the honest statement---but growing far more slowly than
the world it covers.

\paragraph{The stores cross over.} At 25 states the situation store is
\emph{larger} ($8.51$ vs $6.77$). At 144 states it is \emph{smaller} ($14.43$
vs $17.24$) while covering seven times as many situations per item. A
representation that starts out costlier and ends up cheaper, while covering
more, is the signature the measure was built to detect.

\paragraph{Transfer follows coverage.} Situation-keyed memory helps at every
size against both controls, including the shuffled one. Location-keyed memory
helps at no size.

\section{The negative half}

The sharper result is not that location keys fail to help. It is that they
\emph{harm}, and that they are worse than randomly shuffled contents.

\begin{quote}
A shuffled store is harmless \emph{because} it is wrong at random: its keys
rarely match the current state, so it rarely fires. The location store fires
constantly---it was learned from the same \emph{form} of experience it is now
consulted about---and the world behind those matches has changed.
\end{quote}

The obstacles it learned to route around are elsewhere now, so its guidance is
systematically, not randomly, wrong. It steers the agent around obstacles that
are absent and into ones that are present.

Two things follow.

\paragraph{A memory's capacity for harm scales with its coverage.} High-coverage
memory is high-leverage in both directions. This is uncomfortable for the
obvious reason: coverage is also the thing that makes memory worth having.

\paragraph{Retrieval frequency is not evidence of value.} It is highest exactly
where the representation is most misleading---in a shifted environment, which is
the regime memory exists to serve. Reporting hit rate as a virtue inverts the
relationship in the case that matters.

The generalisation we would defend is a design criterion rather than a law:
\emph{index on whatever is invariant in the domain}. Cell coordinates are not
invariant when obstacles move; local geometry is. In a reflexive setting---one
where acting changes the environment---what stays fixed is rarely the state and
more often the relation, which suggests relational indexing is the default
rather than the refinement.

\section{Relation to the companion result}

The same data appear in a companion paper as one of five demonstrations that a
component's \emph{activity} is not evidence of its \emph{contribution}. The
overlap is deliberate and we state it plainly: there, the location store is an
instance of a mechanism that fires and does harm; here, the contrast between
the two keys is used to make generality measurable. The experiment is one
study; the two papers ask different questions of it.

\section{Limits}

\paragraph{One task family.} A grid navigator with random-utility action
selection. The measure is defined for any system with an explicit retrievable
store, but it has been applied once.

\paragraph{Small effects.} The transfer differences are a few percentage points
of success rate. The intervals exclude zero because $n$ is large, not because
the effects are big. The \emph{coverage} result is the striking one---a factor
of seven at the largest size---and it is a property of the representation rather
than of the outcome.

\paragraph{Sub-linear, not invariant.} We predicted the situation store would
stay roughly flat. It grew $1.7\times$. The prediction was directionally right
and quantitatively wrong, and we record it as such.

\paragraph{The measure needs a store.} Coverage per item is defined only for
systems with explicit, countable, retrievable entries. It does not apply to a
representation distributed across weights, which is the case one would most
like to measure. Whether an analogue exists there is open, and we think it is
the interesting question this work does not answer.

\section{Conclusion}

Breadth is coverage bought by accumulation; generality is coverage bought by
structure. If generality is compression, the difference should appear as
description length holding steady while the world grows---and it does, as
coverage per stored item moving in opposite directions with scale.

The result also carries a warning that the positive framing hides. The failing
representation was not inert. It retrieved on nearly every step, and its
retrieving is precisely what made it harmful. Between a representation that
covers many cases and one that covers many \emph{situations} lies the whole
distance between a store that transfers and a store that misleads, and from the
inside---at training time, in the environment where it was learned---the two are
indistinguishable.

\paragraph{Reproducibility.} The preregistration, lock, runner, and result file
are in the accompanying repository. The lock hashes both the preregistration
and the analysis code, and the runner verifies those hashes at startup and
aborts if either has changed.

\end{document}
